\chapter{Brownian Motion}\label{brownian-motion}
\chapterstatus{planned}

The central object of stochastic analysis: its construction, its path properties and its connection
with the Laplace operator.

\noindent\textbf{Prerequisites.} Chapter~\ref{martingales} (Conditional Expectation and Martingales) and Chapter~\ref{general-topology} (General Topology).

\noindent\textbf{On the way to the Hodge conjecture.} This chapter is not needed on the direct
route; it belongs to the reading path of probability
(Subsection~\ref{introduction:subsection-reading-paths}).

\noindent\textbf{Planned contents.}
\begin{enumerate}
\item Gaussian processes; construction of Brownian motion (L\'evy, Kolmogorov's extension and continuity theorems).
\item Path properties: nowhere differentiability, H\"older continuity, quadratic variation, the law of the iterated logarithm.
\item The Markov and strong Markov properties, the reflection principle.
\item Brownian motion and harmonic functions: the Dirichlet problem, recurrence and transience.
\item Donsker's invariance principle.
\item Exercises with solutions.
\end{enumerate}
